\section{Conclusion and future directions}
\label{sec:conclusion}

% From acoustic perception to predictive models of behavioral responsiveness.
% See docs/OUTLINE.md — Conclusion and future directions.
% Draft 1. Style model: Blumstein (agents/STYLE.md).

The evidence reviewed here supports the central claim that acoustic meaning is assigned by the receiver rather than fixed in the waveform, but it supports different parts of that claim with different strength (\tabref{tab:summary-lessons}). Several conclusions are well supported. Naturalness is neither necessary nor sufficient for a sound to affect behavior, and much of what distinguishes call types for laboratory rats is carried by frequency band and gross temporal structure (\secref{sec:q1}). Selective freezing to rat 22-kHz calls depends on prior aversive experience, whereas evidence for unlearned recognition of predator calls in laboratory rodents is inconsistent (\secref{sec:q2}). Neural and physiological measures can separate sounds without any change in behavior, and striatal dopamine responses to aversive sound differ by region and preparation (\secref{sec:q3}). Responses to repeated calls in field sciurids decline in a caller-specific way, and receivers there track caller reliability, although in opposite directions in ground squirrels and marmots (\secref{sec:q5}, \secref{sec:q6}). Other conclusions are weaker: that ecological significance changes the rate of habituation rests mainly on one laboratory study (\secref{sec:q6}), and the effects of irregular timing have been shown only with artificial sounds (\secref{sec:q7}).

A predictive model of responses to repeated sounds would need to specify five components: which sounds are meaningful to which receivers, how a signal's history updates its interpretation, which properties of timing are predictable to the receiver, how responsiveness evolves through repetition and silence, and which laboratory or neural measures predict sustained vigilance or avoidance in the field. The first two are partly characterized. Habituation and identity-specific recovery are well documented in field sciurids, but not against acoustically matched controls; timing and field outcome are largely open. We propose six hypotheses, each designed to discriminate between competing explanations and each pairing field and laboratory systems where possible.

\begin{enumerate}[leftmargin=*, label=(\arabic*)]
  \item \textit{Significance and schedule.} If ecological significance slows habituation or speeds recovery, natural calls should regain their effect after rest more than matched synthetic reconstructions and broadband controls, and this difference should depend on the schedule of presentation (fixed, bounded uniform, or exponential). The decisive measure is the recovery function, not the first response, building on \citet{EvansHammond1983a} and the counterexample of \citet{BerzEtAl2022}.
  \item \textit{Hazard rather than variability.} If receivers habituate partly by predicting when the next sound will occur, then at a matched number of events and mean interval, schedules with constant hazard should preserve responsiveness more than uniformly jittered schedules, which in turn should preserve it more than fixed schedules. Existing studies varied interval variability but not distribution shape \citep{Davis1970,HerryEtAl2007}.
  % VERIFY: WarrenEtAl2025 (not full-text checked): prairie voles discriminate partner from stranger calls.
  \item \textit{Caller reliability in laboratory rodents.} Rats and prairie voles discriminate individual callers \citep{VielleEtAl2021,WarrenEtAl2025}. Manipulating the reliability of individual callers would test whether laboratory rodents discount unreliable callers, as juvenile ground squirrels do \citep{HareAtkins2001}, or scrutinize them, as marmots do \citep{BlumsteinEtAl2004}. If the direction depends on the cost of independent assessment (\secref{sec:q5}), it should change when the opportunity to check for danger is manipulated.
  \item \textit{Regional dopamine responses to ecological calls.} Simultaneous recording from the nucleus accumbens and the tail of the striatum during 22-kHz, 50-kHz, predator, synthetic, and noise stimuli would test whether the opposite regional responses found for noise \citep{FuruyamaEtAl2026} extend to natural calls, and whether accumbens responses to 50-kHz calls \citep{WilluhnEtAl2014} have an aversive counterpart for threat calls.
  % VERIFY: LimaBednekoff1999 (not full-text checked): prediction that effort per high-risk period falls as high-risk periods become frequent.
  \item \textit{Risk allocation with acoustic cues.} Varying the long-term frequency of alarm playback in a field sciurid population, while holding caller identity and acoustic structure constant, would test whether vigilance per alarm falls as alarms become frequent, as the risk allocation hypothesis predicts \citep{LimaBednekoff1999}. We found no such test.
  % VERIFY: LillyEtAl2019 and Mateo2010 (not full-text checked): heterospecific safety cues in gray squirrels; predator-specific physiological responses in Belding's ground squirrels.
  \item \textit{Field computations in laboratory species.} Reliability tracking, eavesdropping, use of heterospecific safety cues, and predator-specific physiological responses are documented in sciurids \citep{HareAtkins2001,GetschowEtAl2013,LillyEtAl2019,Mateo2010}. Reproducing these computations in tractable laboratory species would open them to mechanistic study, and failures to reproduce them would identify where laboratory assays omit the ecological conditions that make a sound meaningful.
\end{enumerate}

These hypotheses also bear on applied acoustic deterrence, in which the persistence of responses to repeated playback determines whether a deterrent remains effective. Comparing timing policies, such as fixed, uniformly jittered, and hazard-based schedules, under matched acoustic exposure would test predictions from hypotheses 1 and 2 directly. No schedule has been shown to prevent habituation to natural calls, and laboratory results with artificial sounds should not be assumed to transfer to free-living animals.
% TODO: disclosure (competing interests) statement, required by the journal, to be added outside this section.

For a receiver, recognizing that a sound indicates danger is the first decision; deciding how long it continues to do so is the harder one. That second decision depends on caller reliability, exposure history, and timing, variables that field studies observe under natural costs and laboratory studies can control and trace to mechanism. Determining whether a sound that once meant danger still does will require both.
